% ---
% id: meat-core/rem:alpha-domain
% kind: rem
% title: "The Rényi order of the key-length statements"
% summary: >-
%   Every key-length statement takes the Rényi order in (1, 2), the domain of
%   the Rényi leftover-hashing lemma, while the entropy, program and gradient
%   results hold on (1, ∞). On (1, 2) the key-length statements and the
%   results they consume apply together.
% prominence: standard
% status: accepted
% uses: [foundations/def:renyi]
% mentions: [thm:keylength, cor:keylength-bound, meat-core/frank-wolfe/cor:fo-keylength, fact:meat]
% assumes: []
% refs:
%   - "kamin2025renyi, Thm. 4, Eq. (13)"
% provenance: newly-derived
% lean:
% ---

\begin{remark}\label{rem:alpha-domain}
Every key-length statement of this contract, among them
Theorem~\ref{thm:keylength}, Corollary~\ref{cor:keylength-bound} and the
key-length corollary of the first-order results
(Corollary~\ref{cor:fo-keylength}), takes $\alpha\in(1,2)$, because the
R\'enyi leftover-hashing lemma their secrecy argument applies is stated
there \cite[Thm.~4, Eq.~(13)]{kamin2025renyi}. The results on entropy
accumulation, the program, its reduction and its gradients are stated on
$(1,\infty)$, where \xref{foundations/def:renyi} defines the sandwiched
entropies, so the intersection of the two domains, on which the key-length
statements and the results they consume apply together, is $(1,2)$; a
result stated on a narrower range of $\alpha$ holds on that range only. The
accumulation import itself extends to $\alpha=\infty$, an endpoint that
\xref{foundations/def:renyi} does not define and no claim uses
(Fact~\ref{fact:meat}, locator note~(vi)).
\end{remark}
